\section{Context Logic}\label{sec:context-logic}

With the help of contextual capabilities algebra, we first define an logic following the standard bunched implication semantics that can help to distinguished the dependent and independent context. We later on extned this logic with expilicit demonic and angelic modalities.

\subsection{Atomic Propositions}

Our context logic is parameterized by a set of atomic propositions $\mathcal{A}$, where we assume there is an corresponding denotation function $\denot{\mathcal{A}}$ maps the atomic proposition $\mathcal{A}$ to set of capabilities that closed under the refinement relation $\sqsubseteq$.
\begin{align*}
  \forall r_1\; r_2. r_1 \in \denot{\mathcal{A}} \land r_2 \sqsubseteq r_1 \implies r_2 \in \denot{\mathcal{A}}
\end{align*}
We highlight that the choice of atomic propositions arbitrary. However, in order to convince for pose examples and further definition of our further context types, we consider a fixed specific atomic proposition domain that is derived from resource-unaware pure logical formula.

\begin{definition}[Pure Qualifier] Assume $\phi$ is a pure logical formula which with respect to environments, i.e., $\sigma(\phi)$ can be either true or false. Then, we define the capability denotation of $\phi$ as the following:
\begin{align*}
  r \in \denot{\phi} \iff \projectTo{r}{\Code{FV}(\phi)} = \{ \sigma ~|~ \Dom{\sigma} = \Code{FV}(\phi) \land \sigma(\phi) \text{ is true} \} 
\end{align*}
Intuitively, a capability $r$ is in the denotation of $\phi$ iff it contains all environments that make $\phi$ true, no missing and no junk.
\end{definition}

\subsection{Basic Context Logic}

\emph{Basic} context logic is parameterized by a set of atomic propositions $\mathcal{\phi}$, where we assume there is an corresponding denotation function $\denot{\phi}$ maps the atomic proposition $\phi$ to set of capabilities that closed under the refinement relation $\sqsubseteq$. We 

Then, every judgment \emph{precisely} describes a resource.
We write $r \models p$ for satisfaction of a resource proposition $p$
at world $r \in \Res$.

\begin{definition}[Basic context logic] Let $(\Res,{+},{\times},\mathbf{1}, \sqsubseteq)$ to be a context algebra. Let $\mathcal{A}$ to be atomic propositions. Let $\denot{-} : \mathcal{A} \sarr \Res$ be an interpretation function of pure formulas $\phi$. Then we have:
\begin{itemize}
\item $r \models \chtrue$ always;
\item $r \models \chfalse$ never;
\item $r \models a$ iff $a \in \mathcal{A}$ and $\denot{a} \sqsubseteq r$;
\item $r \models p \chland q$ iff $r \models p$ and $r \models q$;
\item $r \models p \chlor q$ iff $r \models p$ or $r \models q$;
\item $r \models p \chimpl q$ iff for all $r'$, if $r \sqsubseteq r'$ and $r' \models p$ then $r' \models q$;
% \item $r \models \forall x.p$ iff for all values $v$, $r \models p[x\mapsto v]$;
% \item $r \models \exists x.p$ iff for some $v$, $r \models p[x\mapsto v]$.
\item $r \models p \chstar q$ iff $\exists r_1\;r_2.\; r_1 \times r_2 \sqsubseteq r$ and
  $r_1 \models p$ and $r_2 \models q$.
\item $r \models p \chwand q$ iff $\forall r'.\; \mapCompat{r'}{r}$ and $r' \models p$ implies
  $r' \times r \models q$.
\item $r \models p \choplus q$ iff $\exists r_1\;r_2.\; r_1 + r_2 \sqsubseteq r$ and
  $r_1 \models p$ and $r_2 \models q$.
\item $r \models \forall x.p$ iff $\forall r'.\Dom{r'} = \Dom{r}\cup\{x\} \text{ and }\restrictTo{r'}{\Dom{r}} = r \text{ implies } r' \models p$.
\item $r \models \exists x.p$ iff $\exists r'.\Dom{r'} = \Dom{r}\cup\{x\} \text{ and }\restrictTo{r'}{\Dom{r}} = r \text{ and } r' \models p$.
\end{itemize}
We still write $p \models q$ iff for all $r \in \Res$, $r \models p$ implies $r \models q$.
\end{definition}

Intuitively, $r$ models $\subExt p$ (or $\supExt p$) iff there exists a superset (or subset) $r'$ of $r$ that models $p$.


\begin{definition}[context logic]\label{def:context-logic-modalities} context logic is additionally equipped beyond basic context logic with the overapproximation modality $\subExt$, the underapproximation modality $\supExt$, and the fiberwise universal modality $\forallM{\cdot}$.
  \begin{itemize}
    \item $r \models \subExt p$ iff $\exists r'.\; r \subseteq r'$ and
      $r' \models p$;
    \item $r \models \supExt p$ iff $\exists r'.\; r' \subseteq r$ and
      $r' \models p$,
    \item $r \models \forallM{x}p$ iff $\forall \sigma_x \in \projectTo{r}{\{x\}}. \fiberOver{r}{\sigma_x} \models \sigma_x(p)$;
    \item $r \models \chpers p$ iff $|r| = 1$ and $r \models p$.
    \end{itemize}
\end{definition}

Intuitively, $\forallM{x}p$ is a fiberwise universal modality along the projection
$r \mapsto \projectTo{r}{\{x\}}$.
For each admissible assignment $\sigma_x$ to variable $x$, we look at the
corresponding fiber $\fiberOver{r}{\sigma_x}$ and require that $p$ holds after
substituting $\sigma_x$.
The notation $\forallM{\overline{x_i}} p$ is shorthand for
$\forallM{x_1}\,\forallM{x_2}\,\cdots\,\forallM{x_n}\,p$. The $\forallM{x}p$ means that $p$ holds is independent from the decision made by context resource $r$ on variable $x$.
%  The persistent modality $\chpers p$ means that $p$ holds is independent from the decision made by resource $r$.

% \paragraph{Persistency and disjointness in logic} The persistent modality ($\chpers \phi$) is a stronger modality than disjointness ($\phi' \chstar \phi$); the former requires $\phi$ is independent from any choice made by resource $r$, while the latter only requires a locally independent choice (i.e., independent from $\phi'$). Typically only singleton underapproximation and overapproximation can has persistent modality.


\begin{example}[Dependent Underapproximation] There are three different underapproximations:
  \begin{enumerate}
    \item (dependent) $\supExt (0 < x) \chland \supExt (y < 3)$
    \item (independent, square) $\supExt (0 < x) \chstar \supExt (y < 3)$
    \item (dependent, triangle) $\supExt (0 < x) \chland \supExt (x \sqsubseteq y < 3)\judgfail$. The actual presentation is:
    \begin{align*}
      \supExt (0 < x) \chland \forallM{x}\supExt (x \sqsubseteq y < 3)
    \end{align*}
  \end{enumerate}
We have:
\begin{align*}
  \{[x=1; y=1], [x=1; y=2]\} &\models \forallM{x}\supExt (x \sqsubseteq y < 3)
  \\\{[x=2; y=2]\} &\models \forallM{x}\supExt (x \sqsubseteq y < 3)
  \\\{[x=1; y=1], [x=2; y=2]\} &\models \forallM{x}\supExt (x \sqsubseteq y < 3) \judgfail
  \\&\text{For each fixed-$x$ fiber,}
  \\&\text{all corresponding $y$ outcomes must be present.}
  \\\{[x=1; y=1], [x=1; y=2], [x=2; y=2]\} &\models \forallM{x}\supExt (x \sqsubseteq y < 3) \judgok
  \\&\text{Different $x$-fibers can coexist in the same resource.}
\end{align*}
The $\forallM{x}(0 < x)$ works as a fiberwise universal constraint on $x$: for each
admissible assignment of $x$, we check the corresponding fiber where $x$ is fixed.
Note that it is different from
\begin{align*}
  &r \models \subExt (0 < x) \chland \supExt (x \sqsubseteq y < 3) 
  \\\iff& r \models \subExt (0 < x) \text{ and } r \models \supExt (x \sqsubseteq y < 3)
  \\\implies& \forall \sigma \in r. 0 < \sigma(x) \text{ and } r \models \supExt (x \sqsubseteq y < 3)
  \\\implies& \forall \sigma \in r. 0 < \sigma(x) \text{ and } [x=-1;y=2] \in r \quad \judgfail
\end{align*}\noindent
The problem is that the $\chland$ cannot make the constraint $0 < x$ be added into the pure proposition $x \sqsubseteq y < 3$ (i.e., to be $\supExt (0 < x \land x \sqsubseteq y < 3)$).
\end{example}

% When the expressiveness of atomic propositions is rich enough, nested uses of $\subExt$ and $\supExt$ over logical connectives collapse to $\subExt$ and $\supExt$ applied directly to the underlying pure propositions. For example, $\subExt (\subExt \phi_1 \chland \subExt \phi_2)$ is equivalent to $\subExt \phi_1 \chland \subExt \phi_2$, and $\supExt (\subExt \phi_1 \chland \subExt \phi_2)$ is equivalent to $\supExt (\phi_1 \cap \phi_2)$.

% \begin{theorem}[Collapse of nested modalities] Assume $\mathcal{A}$ is rich enough, i.e.,
% \begin{align*}
%   % &\forall X. \exists \phi \in \mathcal{A}. \denot{\phi} = \{ \subs \mid \Dom{\subs} = X \}
%   % \\
%   &\forall \phi_1\;\phi_2 \in \mathcal{A}. \exists \phi_3 \in \mathcal{A}. \denot{\phi_3} = \denot{\phi_1} \cap \denot{\phi_2}
%   \\& \forall \phi_1\;\phi_2 \in \mathcal{A}. \exists \phi_3 \in \mathcal{A}. \denot{\phi_3} = \denot{\phi_1} \cup \denot{\phi_2}
%   \\& \forall \phi_1\;\phi_2 \in \mathcal{A}. \exists \phi_3 \in \mathcal{A}. \denot{\phi_3} = \denot{\phi_1} \times \denot{\phi_2}
% \end{align*}
% Then, all formulas in context logic can be expressed in approximation normal form, i.e., $\subExt$ and $\supExt$ are applied only to pure propositions.
% \end{theorem}

\subsection{Core Language and context Types}

\begin{figure}[h]
{\small
  \begin{alignat*}{2}
    \text{\textbf{Variables }}& \quad &\quad& x, y, z, \vnu, ...
    \\\text{\textbf{Data constructors }}& \quad &d ::= \quad & () ~|~ \Code{true} ~|~ \Code{false} ~|~ \Code{O} ~|~ \Code{S} ~|~ \Code{Cons} ~|~ \Code{Nil} ~|~ \Code{Leaf} ~|~ \Code{Node}
    \\\text{\textbf{Base Types}}& \quad & b  ::= \quad &  \Code{unit} ~|~ \Code{bool} ~|~ \Code{nat} ~|~ \Code{int} ~|~ ...
    \\\text{\textbf{Basic Types}}& \quad & s  ::= \quad &  b ~|~ s\sarr s
    \\\text{\textbf{Operations}}& \quad & \primop ::= \quad & d ~|~ {+} ~|~ {-} ~|~ {==} ~|~ {<} ~|~ {\sqsubseteq} ~|~ ...
    \\ \text{\textbf{Constants }}& \quad &c ::= \quad & () ~|~ \mathbb{B} ~|~ \mathbb{Z} ~|~ d\ \overline{c}\ ...
    \\ \text{\textbf{Values }}& \quad  & v ::= \quad & c ~|~ x ~|~ \zlam{x}{s}{e} ~|~ \zfix{f}{s}{x}{s}{e}
    \\ \text{\textbf{Expressions}}& \quad & e ::=\quad & v ~|~ \zlet{x}{e}{e} ~|~ \primop\ \overline{v} ~|~  v\;v ~|~ \match{v} \overline{d\ \overline{y} \to e} 
    % \\ & & \quad & 
    % \err ~|~ 
  \end{alignat*}}
\caption{Core Language Syntax. }
\label{fig:core-lang-syntax}
\end{figure}



\begin{figure}[h]
  {
    \begin{alignat*}{2}
      \text{\textbf{Qualifiers}}& \quad & \phi ::= \quad & c ~|~ x ~|~ \primop\;\overline{v} ~|~ \bot ~|~ \top  ~|~ \neg \phi ~|~ \phi \land \phi ~|~ \phi \lor \phi ~|~ \phi \times \phi ~|~ \phi \impl \phi ~|~ \forall x{:}b. \phi ~|~ \exists x{:}b. \phi
      \\ \text{\textbf{context Types}}& \quad & \tau ::= \quad & \ort{b}{\phi} ~|~ \urt{b}{\phi} ~|~ \tau \ctinter \tau ~|~ \tau \ctunion \tau ~|~ \tau \ctoplus \tau ~|~ x{:}\tau \ctsarr \tau ~|~ x{:}\tau \ctwand \tau
      %  ~|~  \ctind \tau 
      ~|~ \ctpers \tau 
      % ~|~ \ctDiv{s}
      \\ \text{\textbf{Type Contexts}}& \quad & \Gamma ::= \quad & \emptyset ~|~ x{:}\tau ~|~ \Gamma \ctstar \Gamma ~|~ \Gamma \ctcomma \Gamma ~|~ \Gamma \ctoplus \Gamma
      \\ \text{\textbf{Contexts With Holes}}& \quad & \Delta ::= \quad & \Box ~|~ \Gamma ~|~ \Delta \ctstar \Delta ~|~ \Delta \ctcomma \Delta ~|~ \Delta \ctoplus \Delta
    \end{alignat*}}
  \vspace*{-.1in}
  {\small
  \begin{align*}
    \text{\textbf{Type erasure}} \quad & \eraserf{\ort{b}{\phi}} \doteq b \quad \eraserf{\urt{b}{\phi}} \doteq b \quad \eraserf{\tau_1 \ctinter \tau_2} \doteq \eraserf{\tau_1} \quad \eraserf{\tau_1 \ctunion \tau_2} \doteq \eraserf{\tau_1}
    \\&\eraserf{\tau_1 \ctoplus \tau_2} \doteq \eraserf{\tau_1} \quad  \eraserf{x{:}\tau_x \ctsarr \tau} \doteq \eraserf{\tau_x} \sarr \eraserf{\tau} \quad  
    % \eraserf{\ctind \tau} \doteq \eraserf{\tau} 
    % \quad \eraserf{\ctpers \tau} \doteq \eraserf{\tau}
    % \quad \eraserf{\ctDiv{s}} \doteq s
    \\& \eraserf{\emptyset} \doteq \emptyset \quad \eraserf{x{:}\tau} \doteq x{:}\eraserf{\tau} \quad \eraserf{\Gamma_1 \ctcomma \Gamma_2} \doteq \eraserf{\Gamma_1}, \eraserf{\Gamma_2} \quad \eraserf{\Gamma_1 \ctstar \Gamma_2} \doteq \eraserf{\Gamma_1}, \eraserf{\Gamma_2}  \quad \eraserf{\Gamma_1 \ctoplus \Gamma_2} \doteq \eraserf{\Gamma_1}
    \\\text{\textbf{Type lifting}} \quad & \liftrf{b} \doteq \ort{b}{\phi} \quad \liftrf{s_x \sarr s} \doteq x{:}\liftrf{s_x} \ctsarr \liftrf{s}  
    \\&\liftrf{\emptyset} \doteq \emptyset \quad \liftrf{x{:}s} \doteq x{:}\liftrf{s} \quad \liftrf{\Gamma_1 , \Gamma_2} \doteq \liftrf{\Gamma_1} \ctcomma \liftrf{\Gamma_2}
    \\\text{\textbf{Type alias}}  \quad & s \doteq \liftrf{s} \quad \tau_{\ctAff} \doteq \tau \ctoplus \liftrf{ \eraserf{\tau} } \quad \prt{b}{\phi} \doteq \ort{b}{\phi} \interty \urt{b}{\phi}
    \\& \tau_x \ctsarr \tau \doteq x{:}\tau_x \ctsarr \tau \quad (\text{when $x$ is not referenced in $\tau$})
    % \quad \tau^{\ctDivM} \doteq \tau \ctoplus \ctDiv{\eraserf{\tau}} 
  \end{align*}}
  \caption{context Types Syntax. }
  \label{fig:context-types-syntax}
  \end{figure}

Here $\ort{b}{\phi}$ is an overapproximate refinement and
$\urt{b}{\phi}$ underapproximate; $\ctoplus$ on types mirrors resource
sum; $\ctsarr$ and $\ctwand$ distinguish ordinary implication from separating
implication at function types.
The context formers $\Gamma_1 \ctstar \Gamma_2$ and
$\Gamma_1 \ctcomma \Gamma_2$ correspond to multiplicative and additive
bunching, respectively. The subtyping and typing judgments are written as
$\Gamma \ctvdash e : \tau$ and $\Gamma \ctvdash \tau_1 <: \tau_2$ in a single declarative system. We omit the well-formedness rules for the context in all typing rules for brevity.

% \paragraph{Persistency and disjointness in type} Again, the persistent modality ($x{:}\ctpers \tau_x$) is a stronger modality than disjointness ($y{:}\tau_y \ctstar x{:}\tau_x$); the former requires judgement $x : \tau_x$ is independent from context, while the latter only requires a locally independent choice (i.e., independent from $y{:}\tau_y$). Here is a example
% \begin{align*}
%   &\Gamma_1 \doteq x{:}\urt{\Int}{\top} \ctcomma (y{:}\urt{\Int}{\top} \ctstar z{:}\urt{\Int}{\top})
%   \\&\Gamma_2 \doteq x{:}\urt{\Int}{\top} \ctcomma (y{:}\urt{\Int}{\top} \ctstar z{:}\ctpers \urt{\Int}{\top})
% \end{align*}
% Consider a program:
% \begin{minted}[fontsize = \footnotesize, escapeinside=??, xleftmargin=18pt, linenos]{ocaml}
%   let x = int_gen () in
%   let y = int_gen () in
%   let z = x in
%   e
% \end{minted}
% when we type check $e$, we can have type context $\Gamma_1$ instead of $\Gamma_2$, since the $z$ is actually dependent on the context (the decision of $x$).

\subsection{Denotational semantics}

\begin{definition}[Totality under a context resource]\label{def:total-under-resource}
For $r \in \Res$ and $e$ with $\FV(e) \subseteq \Dom{r}$, lift $\Code{Total}(e)$ to context logic by:
\[
  r \models \denot{\Code{Total}(e)} \;\iff\;
  \forall \sigma \in r.\; \exists v.\; \mstep{\sigma(e)}{v},
\]
\end{definition}

\begin{align*}
  \\\denot{e}_x &: \Code{Prop_{Strat}} \quad (\text{where $x$ is a free variable})
  \\\denot{e}_x &\doteq \Code{Prop}(\{r \mid \forall \sigma \in r. \mstep{\sigma(e)}{\sigma(x)} \})
  % \\\denot{e}_x &\doteq \forallM{\Code{FV}(e)} \Code{Atom}(\mstep{e}{x})
  \\\denot{\tau} &: \mathcal{E} \sarr \Code{Prop_{Strat}}
  \\\denot{\ort{b}{\phi}} &\doteq \lambda e. \forall \vnu. \denot{e}_\vnu \chimpl \forallM{\Code{FV}(\ort{b}{\phi})}\,\subExt \phi
  \\\denot{\urt{b}{\phi}} &\doteq \lambda e. \forall \vnu. \denot{e}_\vnu \chimpl \forallM{\Code{FV}(\urt{b}{\phi})}\,\supExt \phi
  % \\\denot{x{:}\tau_x \ctsarr \tau} &\doteq \lambda e. \forall y. \denot{e}_y \chimpl \forallM{y} (\denot{\tau_x}\;x \chimpl \denot{\tau}\;(y\;x))
  % \\\denot{x{:}\tau_x \ctwand \tau} &\doteq \lambda e. \forall y. \denot{e}_y \chimpl \forallM{y} (\denot{\tau_x}\;x \chwand \denot{\tau}\;(y\;x))
  \\\denot{x{:}\tau_x \ctsarr \tau} &\doteq \lambda e. \forall x.\denot{\tau_x}\;x \chimpl \denot{\tau}\;(e\;x)
  \\\denot{x{:}\tau_x \ctwand \tau} &\doteq \lambda e. \forall x.\denot{\tau_x}\;x \chwand \denot{\tau}\;(e\;x)
  % \\\denot{\ctind \tau} &\doteq \lambda e. \forallM{\Code{FV}(\tau) \cup \Code{FV}(e)} (\denot{\tau}\;e)
  \\\denot{\ctpers \tau} &\doteq \lambda e. \forall x.\denot{e}_x \chimpl \chpers  (\denot{\tau}\;x)
  \\\denot{\tau_1 \ctoplus \tau_2} &\doteq \lambda e. \denot{\tau_1}\;e \choplus \denot{\tau_2}\;e
  \\\denot{\tau_1 \ctinter \tau_2} &\doteq \lambda e. \denot{\tau_1}\;e \chland \denot{\tau_2}\;e
  \\\denot{\tau_1 \ctunion \tau_2} &\doteq \lambda e. \denot{\tau_1}\;e \chlor \denot{\tau_2}\;e
  \\\denot{\emptyset} &\doteq \Code{Atom}(\textbf{1})
  \\\denot{x{:}\tau} &\doteq \denot{\tau}\;x
  \\\denot{\Gamma_1 \ctcomma \Gamma_2} &\doteq \denot{\Gamma_1} \chland \denot{\Gamma_2}
  \\\denot{\Gamma_1 \ctstar \Gamma_2} &\doteq \denot{\Gamma_1} \chstar \denot{\Gamma_2}
  \\\denot{\Gamma_1 \ctoplus \Gamma_2} &\doteq \denot{\Gamma_1} \choplus \denot{\Gamma_2}
\end{align*}\noindent

\subsection{Declarative typing rules}

\newpage
\section{Basic Typing and Well-Formedness Rules}\label{sec:tech:basic-typing}

\begin{figure}[ht!]
  {\small
  \begin{align*}
    \text{\textbf{Type erasure}} \quad & \eraserf{\ort{b}{\phi}} \doteq b \quad \eraserf{\urt{b}{\phi}} \doteq b \quad \eraserf{\tau_1 \ctinter \tau_2} \doteq \eraserf{\tau_1}
    \\&\eraserf{\tau_1 \ctoplus \tau_2} \doteq \eraserf{\tau_1} \quad  \eraserf{x{:}\tau_x \ctsarr \tau} \doteq \eraserf{\tau_x} \sarr \eraserf{\tau} \quad
    \eraserf{x{:}\tau_x \ctwand \tau} \doteq \eraserf{\tau_x} \sarr \eraserf{\tau} \quad
    \eraserf{\ctpers \tau} \doteq \eraserf{\tau}
    \\& \eraserf{\emptyset} \doteq \emptyset \quad \eraserf{x{:}\tau} \doteq x{:}\eraserf{\tau} \quad \eraserf{\Gamma_1 \ctcomma \Gamma_2} \doteq \eraserf{\Gamma_1}, \eraserf{\Gamma_2} \quad \eraserf{\Gamma_1 \ctstar \Gamma_2} \doteq \eraserf{\Gamma_1}, \eraserf{\Gamma_2}  \quad \eraserf{\Gamma_1 \ctoplus \Gamma_2} \doteq \eraserf{\Gamma_1}
 \end{align*}}
\caption{Type and Typing Context Erasure}
\label{fig:type-erasure}
\end{figure}

\begin{figure}[ht!]
    {\footnotesize
    {\small
   \begin{flalign*}
    &\text{\textbf{Basic Typing }} & \fbox{$\Sigma \basicvdash e : s$}
   \end{flalign*}}
   \begin{prooftree}
   \hypo{}
   \infer1[\textsc{BtConst}]{
   \Sigma \basicvdash c : \Code{BasicTy}(c)
   }
   \end{prooftree}
   \quad
   \begin{prooftree}
   \hypo{ \Sigma(x) = s }
   \infer1[\textsc{BtVar}]{
   \Sigma \basicvdash x : s
   }
   \end{prooftree}
   \quad
   \begin{prooftree}
   \hypo{\Sigma, x{:}s_x \basicvdash e : s}
   \infer1[\textsc{BtFun}]{
   \Sigma \basicvdash \lambda x{:}s_x.e : s_x{\sarr}s
   }
   \end{prooftree}
   \\[.8em]
   \begin{prooftree}
   \hypo{\Sigma, f{:}b\sarr s \basicvdash \lambda x{:}b.e : b\sarr s}
   \infer1[\textsc{BtFix}]{
   \Sigma \basicvdash \zfix{f}{b\sarr s}{x}{b}{e}
   : b{\sarr}s
   }
   \end{prooftree}
   \quad
   \begin{prooftree}
   \hypo{\Code{BasicTy}(\primop) = \overline{s_i}{\sarr}s \quad \forall i. \Sigma \basicvdash v_i : s_i}
   \infer1[\textsc{BtPrimOp}]{
   \Sigma \basicvdash \primop\, \overline{v_i} : s
   }
   \end{prooftree}
   \\[.8em]
   \begin{prooftree}
   \hypo{\Sigma \basicvdash v_1 : s_2 {\sarr} s \quad \Sigma \basicvdash v_2 : s_2}
   \infer1[\textsc{BtApp}]{
   \Sigma \basicvdash v_1\,v_2 : s
   }
   \end{prooftree}
   \\[.8em]
   \begin{prooftree}
   \hypo{\emptyset \basicvdash e_1 : s_x \quad \Sigma, x{:}s_x \basicvdash e_2 : s}
   \infer1[\textsc{BtLetE}]{
   \Sigma \basicvdash \zlet{x:s_x}{e_1}{e_2} : s
   }
   \end{prooftree}
   \quad
   \begin{prooftree}
   \hypo{\Sigma \basicvdash v : s_v \quad
   \forall i, \Code{BasicTy}(d_i) = \overline{s_j} {\sarr} s_v \quad
   \Sigma, \overline{y_j{:}s_j} \basicvdash e_i : s}
   \infer1[\textsc{BtMatch}]{
   \Sigma \basicvdash \match{v} \overline{d_i\,\overline{y_j} \to e_i} : s
   }
   \end{prooftree}
   }
   \caption{Basic Typing Rules}
       \label{fig:basic-type-rules}
   \end{figure}

   \begin{figure}[ht!]
    {\footnotesize
    {\small
   \begin{flalign*}
    &\text{\textbf{Basic Qualifier Typing }} & \fbox{$\Sigma \basicvdash \phi: s$}
   \end{flalign*}}
   \begin{prooftree}
   \hypo{\Code{BasicTy}(c) = s }
   \infer1[\textsc{BtConst}]{
   \Sigma \basicvdash c : s
   }
   \end{prooftree}
   \quad
   \begin{prooftree}
   \hypo{ \Sigma(x) = s }
   \infer1[\textsc{BtVar}]{
   \Sigma \basicvdash x : s
   }
   \end{prooftree}
   \quad
   \begin{prooftree}
   \hypo{ }
   \infer1[\textsc{BtTop}]{
   \Sigma \basicvdash \top : \Code{bool}
   }
   \end{prooftree}
   \quad
   \begin{prooftree}
   \hypo{ }
   \infer1[\textsc{BtBot}]{
   \Sigma \basicvdash \bot : \Code{bool}
   }
   \end{prooftree}
   \\[.8em]
   \begin{prooftree}
   \hypo{\Code{BasicTy}(\primop) = \overline{s_i}{\sarr}s \quad \forall i. \Sigma \basicvdash \phi_i : s_i }
   \infer1[\textsc{BtOp}]{
   \Sigma \basicvdash \primop\,\overline{\phi_i} : s
   }
   \end{prooftree}
   \quad
   \begin{prooftree}
   \hypo{
   \Sigma \basicvdash \phi : \Code{bool} }
   \infer1[\textsc{BtNeg}]{
   \Sigma \basicvdash \neg \phi : \Code{bool}
   }
   \end{prooftree}
   \\[.8em]
   \begin{prooftree}
   \hypo{
   \Sigma \basicvdash \phi_1 : \Code{bool} \quad
   \Sigma \basicvdash \phi_2 : \Code{bool} }
   \infer1[\textsc{BtAnd}]{
   \Sigma \basicvdash \phi_1 \land \phi_2 : \Code{bool}
   }
   \end{prooftree}
   \quad
   \begin{prooftree}
   \hypo{
   \Sigma \basicvdash \phi_1 : \Code{bool} \quad
   \Sigma \basicvdash \phi_2 : \Code{bool} }
   \infer1[\textsc{BtOr}]{
   \Sigma \basicvdash \phi_1 \lor \phi_2 : \Code{bool}
   }
   \end{prooftree}
   \quad
   \begin{prooftree}
   \hypo{
   \Sigma, x{:}b \basicvdash \phi : \Code{bool} }
   \infer1[\textsc{BtForall}]{
   \Sigma \basicvdash \forall x{:}b. \phi : \Code{bool}
   }
   \end{prooftree}
   }
   \caption{Basic Qualifier Typing Rules}
       \label{fig:basic-qualifier-type-rules}
   \end{figure}


\begin{figure}[ht!]
    {\footnotesize
    {\small
   \begin{flalign*}
    &\text{\textbf{Well-Formedness }} & \fbox{$\Sigma \wfvdash \tau$ \quad $\Sigma \wfvdash \Gamma$}
   \end{flalign*}}
   \begin{prooftree}
   \hypo{\Sigma, \vnu{:}b \basicvdash \phi : \Bool}
   \infer1[\textsc{WfBaseOver}]{
    \Sigma \wfvdash \ort{b}{\phi}
   }
   \end{prooftree}
   \quad
   \begin{prooftree}
    \hypo{\Sigma, \vnu{:}b \basicvdash \phi : \Bool}
    \infer1[\textsc{WfBaseUnder}]{
     \Sigma \wfvdash \urt{b}{\phi}
    }
    \end{prooftree}
    \\[.8em]
    \begin{prooftree}
        \hypo{\Sigma \wfvdash \tau_x \quad \Sigma, x{:}\eraserf{\tau_x} \wfvdash \tau}
        \infer1[\textsc{WfArr}]{
         \Sigma \wfvdash x{:}\tau_x \ctsarr \tau
        }
    \end{prooftree}
        \quad
    \begin{prooftree}
        \hypo{\Sigma \wfvdash \tau_x \quad \Sigma \wfvdash \tau}
        \infer1[\textsc{WfWand}]{
            \Sigma \wfvdash x{:}\tau_x \ctwand \tau
        }
    \end{prooftree}
    \quad
    \begin{prooftree}
        \hypo{\Sigma \wfvdash \tau}
        \infer1[\textsc{WfPers}]{
            \Sigma \wfvdash \ctpers \tau
        }
    \end{prooftree}
    \\[.8em]
    \begin{prooftree}
        \hypo{\Sigma \wfvdash \tau_1 \quad \Sigma \wfvdash \tau_2 \quad \eraserf{\tau_1} = \eraserf{\tau_2}}
        \infer1[\textsc{WfInter}]{
         \Sigma \wfvdash \tau_1 \interty \tau_2
        }
    \end{prooftree}
    \quad
    \begin{prooftree}
        \hypo{\Sigma \wfvdash \tau_1 \quad \Sigma \wfvdash \tau_2 \quad \eraserf{\tau_1} = \eraserf{\tau_2}}
        \infer1[\textsc{WfSum}]{
         \Sigma \wfvdash \tau_1 \ctoplus \tau_2
        }
    \end{prooftree}
    \\[.8em]
    \begin{prooftree}
        \hypo{}
        \infer1[\textsc{WfEmp}]{
            \Sigma \wfvdash \emptyset
        }
    \end{prooftree}
    \quad
        \begin{prooftree}
            \hypo{\Sigma \wfvdash \tau_x \quad x \notin \Dom{\Sigma}}
            \infer1[\textsc{WfSingle}]{
             \Sigma \wfvdash x{:}\tau_x
            }
        \end{prooftree}
    \\[.8em]
    \begin{prooftree}
        \hypo{\Sigma \wfvdash \Gamma_1 \quad \Sigma, \eraserf{\Gamma_1} \wfvdash \Gamma_2}
            \infer1[\textsc{WfComma}]{
             \Sigma \wfvdash \Gamma_1 \ctcomma \Gamma_2
        }
    \end{prooftree}
    \quad
    \begin{prooftree}
        \hypo{\Sigma \wfvdash \Gamma_1 \quad \Sigma \wfvdash \Gamma_2}
            \infer1[\textsc{WfStar}]{
             \Sigma \wfvdash \Gamma_1 \ctstar \Gamma_2
        }
    \end{prooftree}
    \\[.8em]
    \begin{prooftree}
        \hypo{\Sigma \wfvdash \Gamma_1 \quad \Sigma \wfvdash \Gamma_2 \quad \eraserf{\Gamma_1} = \eraserf{\Gamma_2}}
            \infer1[\textsc{WfCtxSum}]{
             \Sigma \wfvdash \Gamma_1 \ctoplus \Gamma_2
        }
    \end{prooftree}
   }
   \caption{Well-Formedness Rules for Context Types and Type Contexts}
       \label{fig:basic-refinement-type-rules}
   \end{figure}

   The type erasure function is defined in \autoref{fig:type-erasure}.
   The basic typing rules of our core language, qualifiers and
   refinement types are shown in \autoref{fig:basic-type-rules},
   \autoref{fig:basic-qualifier-type-rules}, and
   \autoref{fig:basic-refinement-type-rules}, resp. We use an
   auxiliary function $\Code{BasicTy}$ to provide a basic type for the
   primitives of our language, e.g., constants, built-in operators,
   and data constructors.


The context formation judgment treats $\ctcomma$ sequentially and $\ctstar$
in parallel: in $\Gamma \basicvdash \Gamma_1 \ctcomma \Gamma_2$, the right
component may refer to bindings introduced by the left component, while in
$\Gamma \basicvdash \Gamma_1 \ctstar \Gamma_2$, both components are checked
only under the ambient $\Gamma$. Thus the two sides of a separating context
cannot refer to variables introduced by each other. We do not give context
holes $\Delta$ an independent well-formedness judgment; rules involving
$\Delta$ are read under the side condition that the plugged contexts are
well formed.

\begin{figure}[h!]
  {\small
  {\normalsize \begin{flalign*}
    &\text{\textbf{Well-formedness rules}} && \fbox{$\Gamma \wfvdash \tau \quad \wfvdash \Gamma$}
  \end{flalign*}}
\mprooftr{28mm}
{$\Gamma \basicvdash \tau$ \quad $\wfvdash \Gamma$}
{WfType}
{$\Gamma \wfvdash \tau$} 
\quad
\mprooftr{28mm}
{$\exists r. \denot{\Gamma}\;r$ \quad $\basicvdash \Gamma$}
{WfCtx}
{$\wfvdash \Gamma$} 
\quad
  {\normalsize \begin{flalign*}
    &\text{\textbf{Subtyping rules}} && \fbox{$\Gamma \ctvdash \tau <: \tau$ \quad $\Gamma \ctsubseteq{X} \Gamma$ }
  \end{flalign*}}
\mprooftr{34mm}
{$\forall e.\; \denot{\Gamma} \chvdash \denot{\tau_1}\;e \chimpl \denot{\tau_2}\;e$}
{Sub}
{$\Gamma \ctvdash \tau_1 <: \tau_2$}
\quad
\mprooftr{36mm}
{$\forall r. r \models \denot{\Gamma_1} \impl \restrictTo{r}{X} \models \denot{\Gamma_2}$ }
{CtxSub}
{$\Gamma_1 \ctsubseteq{X} \Gamma_2$}
% \quad
% 	\mprooftr{22mm}
% 	{$\subExt \denot{\Gamma_1} \models \denot{\Gamma_2}$ }
% 	{CtxOver}
% 	{$\Gamma_1 \ctstar \Gamma_2 = \Gamma_1$}
	  }
	  \caption{Auxiliary well-formedness and semantic subtyping rules for context typing. We disallow the contraction rule (duplicate context) for both $\ctcomma$ and $\ctstar$, since context resources will not be consumed but only depended on. The subtyping rules link typing to semantic implication on denotations. $\Code{ToOver}$ is a context-level over-approximation coercion used to type separating let-bindings.}
  \label{fig:auxiliary-rules}
  \vspace*{-.1in}
\end{figure}

\begin{figure}[h!]
  {\small
  {\normalsize \begin{flalign*}
    &\text{\textbf{Typing rules}} && \fbox{$\Gamma \ctvdash e : \tau$}
  \end{flalign*}}
  % \mprooftr{26mm}
  % {$ \Code{Ty}(\primop) = \tau $}
  % {T-Op}
  % {$\Gamma \ctvdash \primop : \tau$}
  % \quad
  \mprooftr{14mm}
  {$ $}
  {T-Var}
  {$x{:}\tau \ctvdash x : \tau$}
  \quad
  \mprooftr{26mm}
  {$ \emptyset \basicvdash c : b $}
  {T-Const}
  {$\emptyset \ctvdash c : \prt{b}{\vnu = c}$}
  \\[.8em]
  % \mprooftr{34mm}
  % {$\Gamma \ctvdash e_1 : \tau_1 $ \quad $\Delta(\Gamma \ctcomma x{:}\tau_1) \ctvdash e_2 : \tau_2$}
  % {T-Let}
  % {$\Delta(\Gamma) \ctvdash \zlet{x}{e_1}{e_2} : \tau_2$}
  % \quad
  % \mprooftr{40mm}
  % {$\Gamma' \ctvdash e_1 : \tau_1 $ \quad $\Code{ToOver}(\Gamma, \Gamma')$ \quad $\Delta(\Gamma \ctstar x{:}\tau_1) \ctvdash e_2 : \tau_2$}
  % {T-LetD}
  % {$\Delta(\Gamma) \ctvdash \zlet{x}{e_1}{e_2} : \tau_2$}
  % \\[.8em]
  \mprooftr{34mm}
  {$\Gamma \ctvdash e_1 : \tau_1 $ \quad $\Gamma \ctcomma x{:}\tau_1 \ctvdash e_2 : \tau_2$}
  {T-Let}
  {$\Gamma \ctvdash \zlet{x}{e_1}{e_2} : \tau_2$}
  \quad
  \mprooftr{40mm}
  {$\Gamma_1 \ctvdash e_1 : \tau_1 $ \quad $\Gamma_2 \ctstar x{:}\tau_1 \ctvdash e_2 : \tau_2$}
  {T-LetD}
  {$\Gamma_1 \ctstar \Gamma_2 \ctvdash \zlet{x}{e_1}{e_2} : \tau_2$}
  \\[.8em]
  \mprooftr{42mm}
  {$\eraserf{\tau_x} = s_x$ \quad $\Gamma \ctcomma x{:}\tau_x \ctvdash e : \tau$}
  {T-Lam}
  {$\Gamma \ctvdash \zlam{x}{s_x}{e} : x{:}\tau_x \ctsarr \tau$}
  \quad
  \mprooftr{42mm}
  {$\eraserf{\tau_x} = s_x$ \quad  $\Gamma \ctstar x{:}\tau_x \ctvdash e : \tau$}
  {T-LamD}
  {$\Gamma \ctvdash \zlam{x}{s_x}{e} : x{:}\tau_x \ctwand \tau$}
  \\[.8em]
  \mprooftr{40mm}
  {$\Gamma \ctvdash v_1 : x{:}\tau_x \ctsarr \tau$ \quad 
  $\Gamma \ctvdash v_2 : \tau_x$ }
  {T-AppFun}
  {$\Gamma \ctvdash v_1\;v_2 : \tau[x \mapsto v_2]$}
  \quad
  \mprooftr{46mm}
  {$\Gamma_1 \ctvdash v_1 : x{:}\tau_x \ctwand \tau$ \quad 
  $\Gamma_2 \ctvdash v_2 : \tau_x$  }
  {T-AppFunD}
  {$\Gamma_1 \ctstar \Gamma_2 \ctvdash v_1\;v_2 : \tau[x \mapsto v_2]$}
  \\[.8em]  
  \mprooftr{18mm}
  {$\Gamma \ctvdash e : \tau_1$ \quad $\Gamma \ctvdash \tau_1 <: \tau_2$}
  {T-Sub}
  {$\Gamma \ctvdash e : \tau_2$}
  \quad
  \mprooftr{23mm}
  {$\Gamma_2 \ctvdash e : \tau$ \quad $\Gamma_1 \ctsubseteq{\Code{FV}(e) \cup \Code{FV}(\tau)} \Gamma_2$}
  {T-CtxSub}
  {$\Gamma_1 \ctvdash e : \tau$}
  \quad
  \mprooftr{34mm}
  {$\Code{Ty}(\primop) = \overline{x_i{:}\tau_i}\ctsarr \tau_y$ \quad 
    $\forall i. \Gamma \ctvdash v_i : \tau_i$ }
  {T-AppOp}
  {$\Gamma \ctvdash \primop\;\overline{v_i} : \tau[\overline{x_i \mapsto v_i}]$}
  \\[.8em]
  \mprooftr{64mm}
  {$\eraserf{\tau_x} = b$ \quad $\eraserf{\tau} = t$ \quad $\Gamma \ctcomma x{:}\tau_x \ctcomma f{:}(x{:}(\tau_x \ctinter \ort{b}{\vnu \prec x} ) \ctsarr \tau) \ctvdash e : \tau$}
  {T-Fix}
  {$\Gamma \ctvdash \zfix{f}{b\sarr t}{x}{b}{e} : x{:}\tau_x \ctsarr \tau$}
  \\[.8em]
  \mprooftr{64mm}
  {$\eraserf{\tau_x} = b$ \quad $\eraserf{\tau} = t$ \quad  $\Gamma \ctstar (x{:}\tau_x \ctcomma f{:}(x{:}(\tau_x \ctinter \ort{b}{\vnu \prec x} ) \ctsarr \tau)) \ctvdash e : \tau$}
  {T-FixD}
  {$\Gamma \ctvdash \zfix{f}{b\sarr t}{x}{b}{e} : x{:}\tau_x \ctwand \tau$}
  \\[.8em]
  \mprooftr{74mm}
  {$\forall i.\;\Code{Ty}(d_i) = \overline{y{:}b_y}\sarr \tau_i$ and $\Delta_i(\Gamma_i) \ctvdash v : \tau_i$ and $\Delta_i(\Gamma_i) \ctvdash e_i : \tau_i$ }
  {T-Match}
  {$\ctbigoplus \overline{\Gamma_i} \ctvdash \match{v} \overline{d\ \overline{y} \to e}  : \ctbigoplus \overline{\tau_i}$}
  }
  \caption{Declarative typing rules for the core language.}
  \label{fig:declaritive-typing-rules}
  \vspace*{-.1in}
\end{figure}

\begin{figure}[h!]
{\small
{\normalsize
\begin{flalign*}
  &\text{\textbf{Derived Structural rules (optional)}} && \fbox{$\Gamma \equiv \Gamma$ \quad $\Code{Pres}(\tau)$ }
\end{flalign*}}
\mprooftr{22mm}
{$\Gamma_1 \ctvdash e : \tau$ \quad $\Gamma_1 \ctcomma \Gamma_2 \ctsubseteq{\Dom{\Gamma_1}} \Gamma_1$}
{W$_{\ctcomma}$}
{$\Gamma_1 \ctcomma \Gamma_2 \ctvdash e : \tau$}
\quad
\mprooftr{22mm}
{$\Gamma_1 \ctvdash e : \tau$ \quad $\Gamma_1 \ctstar \Gamma_2 \ctsubseteq{\Dom{\Gamma_1}} \Gamma_1$}
{W$_{\ctstar}$}
{$\Gamma_1 \ctstar \Gamma_2 \ctvdash e : \tau$}
\quad
\mprooftr{22mm}
{$\Gamma_2 \ctcomma \Gamma_1 \ctvdash e : \tau$ \quad $\Gamma_1 \ctcomma \Gamma_2 \ctsubseteq{} \Gamma_2 \ctcomma \Gamma_1$}
{E$_{\ctcomma}$}
{$\Gamma_1 \ctcomma \Gamma_2 \ctvdash e : \tau$}
\quad
\mprooftr{22mm}
{$\Gamma_2 \ctstar \Gamma_1 \ctvdash e : \tau$ \quad $\Gamma_1 \ctstar \Gamma_2 \ctsubseteq{} \Gamma_2 \ctstar \Gamma_1$}
{E$_{\ctstar}$}
{$\Gamma_1 \ctstar \Gamma_2 \ctvdash e : \tau$}
\\[.8em]
\mprooftr{22mm}
{$\Gamma_1, \Gamma_2 \ctvdash e : \tau$ \quad $\Gamma_1 \ctsubseteq{\Dom{\Gamma_1}} \Gamma_1 \ctcomma \Gamma_2$}
{C$_{\ctcomma}$}
{$\Gamma_1 \ctvdash e : \tau$}
\quad
\mprooftr{22mm}
{$\Gamma_1 \ctstar \Gamma_2 \ctvdash e : \tau$ \quad $\Gamma_1 \ctsubseteq{\Dom{\Gamma_1}} \Gamma_1 \ctstar \Gamma_2$}
{C$_{\ctstar}$}
{$\Gamma_1 \ctvdash e : \tau$}
\quad
\mprooftr{26mm}
{$\tau = \ort{b}{\phi}$ or $\eraserf{\tau} = s_1 \sarr s_2$}
{Pers}
{$\Code{Pres}(\tau)$}
\quad
\mprooftr{24mm}
{$\Code{Pres}(\tau_x)$}
{PersM}
{$x{:}\tau_x \ctcomma \Gamma \equiv x{:}\tau_x \ctstar \Gamma$}
\\[.8em]
\mprooftr{22mm}
{$ $}
{Split$_{\ctstar}$}
{$\Gamma_1 \ctstar (\Gamma_2 \ctoplus \Gamma_3) \equiv (\Gamma_1 \ctstar \Gamma_2) \ctoplus (\Gamma_1 \ctstar \Gamma_3)$}
\quad
\mprooftr{22mm}
{$\Code{Pres}(\tau)$}
{SplitPers}
{$x{:}\tau \equiv x{:}\tau \ctoplus \tau$}
\quad
\mprooftr{38mm}
{$ $}
{SplitMerge}
{$x{:}\urt{b}{\phi_1 \lor \phi_2} \equiv x{:}\urt{b}{\phi_1} \ctoplus x{:}\urt{b}{\phi_2}$}
\\[.8em]
\mprooftr{22mm}
{$ $}
{Split$_{\ctcomma}$}
{$\Gamma_1 \ctcomma (\Gamma_2 \ctoplus \Gamma_3) \equiv (\Gamma_1 \ctcomma \Gamma_2) \ctoplus (\Gamma_1 \ctcomma \Gamma_3)$}
\quad
\mprooftr{22mm}
{$\Gamma \ctvdash e : \tau$ \quad $\Delta(\Gamma) \ctsubseteq{\Dom{\Gamma}} \Gamma $}
{Frame}
{$\Delta(\Gamma) \ctvdash e : \tau$}
\quad
\mprooftr{35mm}
{$\Gamma \ctcomma x{:}\ort{b}{\phi} \ctvdash e : \tau$}
{Forall}
{$\Gamma \ctvdash e : \Code{Forall}(x{:}\phi, \tau)$}
\\[.8em]
\mprooftr{35mm}
{$\Gamma \ctcomma x{:}\urt{b}{\phi} \ctvdash e : \tau$}
{Exists}
{$\Gamma \ctvdash e : \Code{Exists}(x{:}\phi, \tau)$}
\quad
\mprooftr{65mm}
{$\forall i.\;\Gamma_i \ctvdash v : \tau_{vi}$  and $\Dom{\Delta_i} = \{\overline{y_{ij}}\}$ and  $\Delta_i(\Gamma_i) \ctvdash d_i(\overline{y_{ij}}) : \tau_{vi}$ and $\Delta_i(\Gamma_i) \ctvdash e_i : \ort{b}{\phi}$ }
  {T-MatchO}
  {$\ctbigoplus \overline{\Gamma_i} \ctvdash \match{v} \overline{d_i\ \overline{y_{ij}} \to e_i}  : \ort{b}{\phi}$}
  \\[.8em]
\mprooftr{65mm}
  {$\forall i.\;\Gamma_i \ctvdash v : \tau_{vi}$  and $\Dom{\Delta_i} = \{\overline{y_{ij}}\}$ and  $\Delta_i(\Gamma_i) \ctvdash d_i(\overline{y_{ij}}) : \tau_{vi}$ and $\Delta_i(\Gamma_i) \ctvdash e_i : \urt{b}{\phi_i}$ }
  {T-MatchU}
  {$\ctbigoplus \overline{\Gamma_i} \ctvdash \match{v} \overline{d_i\ \overline{y_{ij}} \to e_i}  : \urt{b}{\bigvee \phi_i}$}
}
	\caption{Derived rules.}
\label{fig:declaritive-typing-rules}
\vspace*{-.1in}
\end{figure}

\begin{theorem}[Well-formed operator context] The operator context $\Phi$ is well-formed iff for all primtive operator $\primop$ we have:
{\begin{align*}
    &\Phi(\primop) = \overline{x{:}\ort{b_x}{\phi_x}}\ctsarr \prt{b}{\phi} \text{ and } \models \denot{\overline{x{:}\ort{b_x}{\phi_x}}} \chiff \denot{\prt{b}{\phi}}\; (\primop\;\overline{x})
  \end{align*}}\noindent
\end{theorem}

\begin{theorem}[Fundamental theorem]
  \begin{align*}
    \Gamma \ctvdash e : \tau \text{ implies } \denot{\Gamma} \chvdash \denot{\tau}\;e \chland \Code{Total}(e)
  \end{align*}\noindent
\end{theorem}

\begin{theorem}[Denotational Soundness]
  \begin{align*}
    \emptyset \ctvdash e : \tau \text{ implies } \forall x. \{[x\mapsto v] ~|~ \mstep{e}{v}\} \models \denot{\tau}\; x
  \end{align*}\noindent
\end{theorem}


\begin{corollary}[Simulate different reasoning styles] The following facts hold:
  \begin{enumerate}
    \item $\emptyset \ctvdash e : \ort{b}{\phi}$ implies that $\forall v. \mstep{e}{v} \impl \phi[\vnu \mapsto v]$.
    \item $\emptyset \ctvdash e : \urt{b}{\phi}$ implies that $\forall v. \phi[\vnu \mapsto v] \impl \mstep{e}{v}$.
    \item Refinement type check: $\emptyset \ctvdash v_f :  x{:}\ort{b_x}{\phi_x} \ctsarr \ort{b}{\phi}$ implies that $\forall v_x\;v. \phi_x[\vnu \mapsto v_x] \land \mstep{v_f\;v_x}{v} \impl \phi[\vnu \mapsto v]$.
    \item Coverage type check: $\emptyset \ctvdash v_f :  x{:}\ort{b_x}{\phi_x} \ctsarr \urt{b}{\phi}$ implies that $\forall v_x\;v. \phi_x[\vnu \mapsto v_x] \land \phi[\vnu \mapsto v]  \impl \mstep{v_f\;v_x}{v} $.
    \item Incorrectness type check: $\emptyset \ctvdash v_f : \urt{b_x}{\phi_x} \ctsarr \urt{b}{\phi}$ implies that \\ $\forall v. \phi[\vnu \mapsto v] \impl \exists v_x. \phi_x[\vnu \mapsto v_x] \land  \mstep{v_f\;v_x}{v} $.
    \item Sufficient incorrectness: $\emptyset \ctvdash v_f : \ort{b_x}{\phi_x} \ctsarr \ort{b}{\phi} \ctoplus b$ implies that \\ $\forall v_x. \phi_x[\vnu \mapsto v_x] \impl \exists v. \phi[\vnu \mapsto v] \land \mstep{v_f\;v_x}{v} $.
    \item Existential type check: $\emptyset \ctvdash v_f : \urt{b_x}{\phi_x} \ctsarr \ort{b}{\phi}\ctoplus b$ implies that \\ $\exists v_x\;v. \phi_x[\vnu \mapsto v_x] \land \phi[\vnu \mapsto v] \land \mstep{v_f\;v_x}{v} $.
  \end{enumerate}
\end{corollary}

\newpage
\subsection{Examples}

\paragraph{Function applying multiple times} Consider the following example:
\begin{minted}[xleftmargin=5pt, numbersep=4pt, linenos = true, fontsize = \footnotesize, escapeinside=??]{OCaml}
    let f = fun () -> 1 ?$\oplus$? 2 in
    let x = f () in
    let y = f () in
    x + y
\end{minted}
We have
\begin{align*}
  &f{:}\Unit\sarr\urt{\Nat}{\vnu = 1 \lor \vnu = 2} \vdash f : \Unit\sarr\urt{\Nat}{\vnu = 1 \lor \vnu = 2}
  \\&\underline{\emptyset \vdash () : \Unit}
  \\&\underline{\emptyset \ctstar f{:}\Unit\sarr\urt{\Nat}{\vnu = 1 \lor \vnu = 2} \vdash f(): \urt{\Nat}{\vnu = 1 \lor \vnu = 2}}
  \\&\underline{f{:}\Unit\sarr\urt{\Nat}{\vnu = 1 \lor \vnu = 2} \vdash f(): \urt{\Nat}{\vnu = 1 \lor \vnu = 2}}
  \\&\underline{\emptyset \ctstar x{:}\urt{\Nat}{\vnu = 1 \lor \vnu = 2} \vdash \zlet{y}{f()}{ x + y}: \urt{\Nat}{\vnu = 1 \lor \vnu = 2}} \quad \judgfail
  \\&\emptyset \ctstar f{:}\Unit\sarr\urt{\Nat}{\vnu = 1 \lor \vnu = 2} \vdash \zlet{x}{f()}{ ...}: \urt{\Nat}{\vnu = 1 \lor \vnu = 2}
\end{align*}\noindent
What we want:
\begin{align*}
  &\underline{f{:}\Unit\sarr\urt{\Nat}{\vnu = 1 \lor \vnu = 2} \ctstar f{:}\Unit\sarr\urt{\Nat}{\vnu = 1 \lor \vnu = 2} \vdash ... : ...} 
  \\&f{:}\Unit\sarr\urt{\Nat}{\vnu = 1 \lor \vnu = 2} \vdash ... : ...
\end{align*}\noindent
Thus, the function $f$ is duplicatable.

\paragraph{What provious $\ctpers$ modality does? How our type system does ?} It means this typing can be duplicate and has no different with $\ctstar$ and $\ctcomma$. Currently, the over-base type and function type are persistent. Although we disallowing duplication, we can have duplicates during spliting. For example, consider:
\begin{align*}
  &\ctvdash (\lambda ().1) \ctoplus (\lambda ().2) : \Unit \sarr \urt{\Int}{\vnu = 1} \ctoplus \Unit \sarr \urt{\Int}{\vnu = 2} \quad \judgok 
  \\&\ctvdash (\lambda ().1) \ctoplus (\lambda ().2) : \Unit \sarr \urt{\Int}{\vnu = 1 \lor \vnu = 2} \quad \judgfail
\end{align*}\noindent

\paragraph{What provious $\ctstar$ modality does? How our type system does ?} It means the type judgement is independent of the context. Currently, we directly using framing rule to indicates the dependency requirement:

\mprooftr{22mm}
  {$ \Gamma \vdash e : \tau$}
  {Frame}
  {$\Gamma \ctstar \Gamma' \vdash e : \tau$}

\paragraph{What we lose without $\ctstar$ modality?} Consider a function $f$:
\begin{align*}
  f \doteq \zlam{x}{\Int}{\Code{int\_gen()}}
\end{align*}\noindent
Then, we have:

\mprooftr{65mm}
  {$ x{:}\urt{\Int}{\top} \vdash f\; x : \urt{\Int}{\top}$ \quad $x{:}\urt{\Int}{\top} \ctcomma y{:}\urt{\Int}{\top} \ctvdash e : \tau$}
  {TApp}
  {$x{:}\urt{\Int}{\top} \ctstar \zlet{y}{f\;x}{e} : \tau$}

Here, the $y$ is always dependent on the arguement $x$, whatever normal arrow or magic wand the type of $f$ is (they can only distinguish the dependency of the parameter $x$).
\begin{align*}
  &\ctvdash f : x{:}\urt{\Int}{\top} \ctsarr \urt{\Int}{\top}
  \\&\ctvdash f : x{:}\urt{\Int}{\top} \ctwand \urt{\Int}{\top}
\end{align*}\noindent

In this example, the assignment of $x$ is indeed independent of $y$, but the type system cannot tell this. Thus, we want:

\mprooftr{65mm}
  {$ x{:}\urt{\Int}{\top} \vdash f\; x : \ctstar \urt{\Int}{\top}$ \quad $x{:}\urt{\Int}{\top} \ctstar y{:}\urt{\Int}{\top} \ctvdash e : \tau$}
  {TApp}
  {$x{:}\urt{\Int}{\top} \ctstar \zlet{y}{f\;x}{e} : \tau$}

Moreover, in order to type check the $\ctstar$ modality, we need to add the following rule:

\mprooftr{65mm}
  {$ x{:}\ort{\Int}{\top} \vdash f\; x : \urt{\Int}{\top}$}
  {TFlip}
  {$ x{:}\urt{\Int}{\top} \vdash f\; x : \ctstar \urt{\Int}{\top}$}

We lose the ability to reason about this kind of terms, although the function can has a more precise type:
\begin{align*}
  \ctvdash f : x{:}\ort{\Int}{\top} \ctsarr \urt{\Int}{\top}
\end{align*}\noindent
And rewrite the function application as:
\begin{align*}
  &x{:}(\ort{\Int}{\top} \ctinter \urt{\Int}{\top}) \ctvdash x : \ort{\Int}{\top} \ctinter \urt{\Int}{\top}
  \\&\underline{\emptyset \ctvdash f : x{:}\ort{\Int}{\top} \ctsarr \urt{\Int}{\top}}
  \\&\underline{x{:}\ort{\Int}{\top} \ctvdash f\;x : \urt{\Int}{\top}}
  \\&\underline{x{:}\ort{\Int}{\top} \ctstar y{:}\urt{\Int}{\top} \ctvdash \zlet{y}{f\;x}{e} : \tau}  
\end{align*}\noindent

\mprooftr{65mm}
  {$ \emptyset \vdash f\; x : \urt{\Int}{\top}$ \quad $x{:}\urt{\Int}{\top} \ctstar \Box (\emptyset \ctcomma y{:}\urt{\Int}{\top}) \ctvdash e : \tau$}
  {TApp}
  {$x{:}\urt{\Int}{\top} \ctstar \Box (\emptyset) \ctvdash \zlet{y}{f\;x}{e} : \tau$}

% \mprooftr{65mm}
%   {$ x{:}\ort{\Int}{\top} \vdash f\; x : \urt{\Int}{\top}$ \quad $x{:}\ort{\Int}{\top} \ctcomma y{:}\urt{\Int}{\top} \ctvdash e : \tau$}
%   {TApp}
%   {$x{:}\ort{\Int}{\top} \ctvdash \zlet{y}{f\;x}{e} : \tau$}


\paragraph{No quotient modality anymore} We flip the underapproximation type $x{:}\urt{b}{\phi}$ directly into $x{:}\ort{b}{\top}$. If users want a tighter safety constraint, they should explicitly state it, e.g.,
\begin{align*}
  x{:}(\urt{\Int}{\vnu = 3} \interty \ort{\Int}{\vnu > 0}) ...
\end{align*}\noindent
after flipping we get
\begin{align*}
  x{:}(\ort{\Int}{\top} \interty \ort{\Int}{\vnu > 0}) ...
\end{align*}\noindent

\paragraph{Existential parameter} We need to distinguish the difference between ``there exists a parameter'' and ``an input parameter that has an underapproximation guarantee (incorrectness logic)''. Consider the following example:
\begin{align*}
  &\vdash f : x{:}\urt{\Int}{\vnu > 0} \sarr \urt{\Int}{\vnu > 0}
  \\\text{\textcolor{blue}{Incorrectness style:} }& \text{if inputs cover ``positive numbers'', the output also covers ``positive numbers''.}
  \\\text{\textcolor{DeepGreen}{Exists style:} }& \text{there exists one ``positive'' input such that the output also covers ``positive numbers''.}
  \\& f_1 = \lambda x.x \quad (\text{Incorrectness}\; \judgok) \quad (\text{Exists}\; \judgfail)
  \\& f_2 = \lambda x.\Code{int\_gen()} \quad (\text{Incorrectness}\; \judgok) \quad (\text{Exists}\; \judgok)
\end{align*} 
We need another way to express ``exists-style''. One proposal is to use $\choplus$:
\begin{align*}
  &\vdash f : x{:}\urt{\Int}{\vnu > 0} \sarr \ort{\Int}{\vnu > 0} \choplus \ort{\Int}{\top} 
  \\\text{\textcolor{DeepGreen}{Exists style:} }& \text{there exists one ``positive'' input such that the output is also ``positive''.}
  \\& \text{The resource is non-empty, thus we cannot also choose the second case in $\choplus$.}
  \\&\vdash f : x{:}\urt{\Int}{\vnu > 0} \sarr \urt{\Int}{\vnu > 0} \choplus \ort{\Int}{\top} 
  \\\text{\textcolor{DeepGreen}{Exists style:} }& \text{exists ``positive'' input makes the output cover ``positive numbers''.}
  \\& f_1 = \lambda x.x \quad (\text{Incorrectness}\; \judgok) \quad (\text{Exists}\; \judgfail)
  \\& f_3 = \lambda x.(x - 1) \quad (\text{Incorrectness}\; \judgok) \quad (\text{Exists}\; \judgok)
  \\& \supExt(x > 0) = \supExt(x > 1) \oplus \supExt(x = 1)  
\end{align*} 
The problem here is that we do not enforce that there exists only \emph{one} input, which can be addressed by the persistency modality (correct under singleton resources):
\begin{align*}
  &\vdash f : x{:}\urt{\Int}{\vnu > 0} \sarr (\chpers \urt{\Int}{\vnu > 0}) \choplus \ort{\Int}{\top} 
  \\\text{\textcolor{DeepGreen}{Exists style:} }& \text{exists ``positive'' input makes the output cover ``positive numbers''.}
  \\& f_3 = \lambda x.(x - 1) \quad (\text{Incorrectness}\; \judgok) \quad (\text{Exists}\; \judgfail)
  \\& f_2 = \lambda x.\Code{int\_gen()} \quad (\text{Incorrectness}\; \judgok) \quad (\text{Exists}\; \judgok)
\end{align*}\noindent
We discuss the details of persistency later.

\paragraph{Persistent modality and persistent proposition} \emph{Persistency} means ``the proof does not rely on any resource ownership (resource consumption)'', so this result can be used multiple times. A \emph{persistent property} means that the property itself implies that it is persistent. Formally, we have:
\begin{align*}
  r \models \chpers p &\text{ iff } \forall \sigma \in r. \{\sigma \} \models p\\
  \Code{Persistent}(p) &\text{ iff } p \models \chpers p 
\end{align*}\noindent
Intuitively, the overapproximation modality is persistent, since it can be duplicated arbitrarily. This fact can be proved as follows:
\begin{align*}
  &\subExt p \models \chpers \subExt p
  \\\iff& \forall r. r \models \subExt p \implies r \models \chpers \subExt p
  \\\iff& \forall r. r \models \subExt p \implies \forall \sigma \in r. \{\sigma\} \models \subExt p
  \\\text{Note that: }& \{\sigma\} \subseteq r
  \\\text{Thus: }& \forall \sigma \in r. \{\sigma\} \models \subExt \subExt p = \subExt p 
\end{align*}\noindent
Note that I have modified the definition of persistent property to be more general, previously:
\begin{align*}
  r \models \chpers p &\text{ iff } |r| = 1 \land r \models p
\end{align*}\noindent
which disallows
\begin{align*}
  \{ [x=1], [x=2] \} \models \chpers \subExt (x > 0)
\end{align*}\noindent
where the right-hand side should obviously be a persistent property. We would like to highlight that the persistent modality is different from a persistent proposition (or type). The over type $\ort{b}{\phi}$ is a persistent type, and the under type $\urt{b}{\phi}$ is not (unless it is a singleton under type). However, we can have $\chpers \urt{b}{\phi}$, which means that the proof of this type judgment does not rely on any resource ownership.

According to the definition of persistent modality, we find that it indeed ``flips'' a resource into its ``for all'' style version, which builds a deep connection between the underapproximation and overapproximation modalities. Precisely, we have:
\begin{align*}
  &\Gamma \vdash e_x : \chpers \tau_x
  \\&\underline{\Gamma \chstar x{:}\tau_x \vdash e : \tau}
  \\&\Gamma \vdash \zlet{x}{e_x}{e} : \tau
\end{align*}\noindent
Since the proof of $\Gamma \vdash e_x : \chpers \tau_x$ does not rely on any resource ownership, it can be made disjoint from the type context. Moreover, we have:
\begin{align*}
  &\underline{\Code{flip}(\Gamma) \vdash e : \tau}
  \\&\Gamma \vdash e : \chpers \tau
\end{align*}\noindent
which is similar to requiring that a proposition hold under all possible singleton resources; here we use the ``overapproximation modality'' to express the ``for all'' style version. As a concrete example, we have:
\begin{align*}
  &\underline{y{:}\ort{\Int}{\top} \vdash \Code{int\_gen()} : \urt{\Int}{\top}} \judgok
  &\underline{y{:}\ort{\Int}{\top} \vdash y : \urt{\Int}{\top}} \judgfail
  \\&y{:}\urt{\Int}{\top} \vdash \Code{int\_gen()} : \chpers \urt{\Int}{\top}
  &y{:}\urt{\Int}{\top} \vdash y : \chpers \urt{\Int}{\top}
  \\&\underline{y{:}\urt{\Int}{\top} \chstar x{:}\urt{\Int}{\top} \vdash e : \tau}
  &\underline{y{:}\urt{\Int}{\top} \chstar x{:}\urt{\Int}{\top} \vdash e : \tau}
  \\&y{:}\urt{\Int}{\top} \vdash \zlet{x}{\Code{int\_gen()}}{e} : \tau
  &y{:}\urt{\Int}{\top} \vdash \zlet{x}{y}{e} : \tau
\end{align*}\noindent
Note that the type of $x$ in the type context cannot be persistent again, so it cannot be duplicated anymore \emph{after binding to a specific variable}. In a call-by-value semantics, a variable cannot be bound to two different values. When can $x$ have a persistent type? Only when $x$ is already bound to a specific value.
\begin{align*}
  &\Gamma \vdash v_x : \chpers \tau_x
  \\&\underline{\Gamma \chstar x{:}\chpers \tau_x \vdash e : \tau}
  \\&\Gamma \vdash \zlet{x}{v_x}{e} : \tau
\end{align*}\noindent
For example,
\begin{align*}
  &\emptyset \vdash \lambda ().\Code{int\_gen()} : \chpers (\Unit\sarr\urt{\Int}{\top})  \judgok
  \\&\underline{x{:}\chpers (\Unit\sarr\urt{\Int}{\top}) \vdash e : \tau}
  \\&\emptyset \vdash \zlet{x}{\lambda ().\Code{int\_gen()}}{e} : \tau
  \\
  \\&\emptyset \vdash \lambda ().\Code{nat\_gen()} \oplus \lambda ().\Code{neg\_gen()} : \chpers (\Unit\sarr\urt{\Int}{\top}) \judgfail
  \\&\underline{x{:}\chpers (\Unit\sarr\urt{\Int}{\top}) \vdash e : \tau}
  \\  &\emptyset \vdash \zlet{x}{\lambda ().\Code{nat\_gen()} \oplus \lambda ().\Code{neg\_gen()}}{e} : \tau
\end{align*}\noindent

\paragraph{Some consideration about persistency} In current setting, only persistency means holds under singleton resource (it is really hard to generalize and somehow counter-intuitive). One straightforward result is ``$\subExt \phi$'' is not persistent. Consider several examples:
\begin{align*}
  &\subExt (x > 0)\chland \supExt (y > 0) \models \subExt (x > 0)\chstar \supExt (y > 0) \judgfail
\end{align*}\noindent
Consider $m \doteq \{ [x=1;y=1], [x=2;y=2], ... \}$ which holds the left-hand side, but not the right-hand side. Actually, in this case $x$ and $y$ are still dependent, if we decide the value of $x$, the value of $y$ will also be determined. In general, 
\begin{align*}
  &\chpers \subExt (x > 0)\chland \supExt (y > 0) \models \chpers\subExt (x > 0)\chstar \supExt (y > 0) \judgok
\end{align*}\noindent
where I call $\chpers \subExt \phi$ as \emph{universal} modality. For example,
\begin{align*}
  &\denot{\subExt (0 \sqsubseteq x \sqsubseteq 1)} =\{ \{ [x=0] \},  \{ [x=1] \}, \{ [x=0], [x=1] \}, ... \}
  \\&\denot{\chpers \subExt (0 \sqsubseteq x \sqsubseteq 1)} =\{ \{ [x=0] \},  \{ [x=1] \}, ... \}
\end{align*}\noindent
which means that the universal modality is persistent. When overapproximation appears at negative position, they are equivalent to the persistent modality in some cases:
\begin{align*}
  &\subExt p \models \subExt q \iff \chpers \subExt p \models \subExt q
  \\&\subExt p \models \supExt q \iff \chpers \supExt p \models \supExt q
\end{align*}\noindent
However, in general, they are not equivalent. For example,
\begin{align*}
  &\chpers \subExt (0 \sqsubseteq x \sqsubseteq 1) \models (x = 0) \chlor (x = 1) \judgok
  \\&\subExt (0 \sqsubseteq x \sqsubseteq 1) \models (x = 0) \chlor (x = 1) \judgfail
\end{align*}\noindent
Here is another typing example:
\begin{align*}
  & \vdash \lambda x.x : x{:}\ort{\Bool}{\top} \sarr \ort{\Bool}{\vnu = \true} \unionty \ort{\Bool}{\vnu = \false} \judgfail
  \\&\vdash \lambda x.x : x{:}(\chpers\ort{\Bool}{\top}) \sarr \ort{\Bool}{\vnu = \true} \unionty \ort{\Bool}{\vnu = \false} \judgok
\end{align*}\noindent
It means that when $\lambda x.x$ apply to $\true$ or $\false$, the result is either $\true$ or $\false$.  However, when $\lambda x.x$ applies to $\Code{bool\_gen()}$, the result is not guaranteed to be either $\true$ or $\false$ (can be both).
This issue comes from the following fact that ``logical disjunction'' is not the same as `Semantical typing-level union'' in an non-deterministic language.
\begin{align*}
  & \denot{\ort{b}{\phi_1 \lor \phi_2}} \not= \denot{\ort{b}{\phi_1} \unionty \ort{b}{\phi_2}} 
\end{align*}\noindent
which is pointed out by Jana Dunfield and Frank Pfenning\cite{UnionTypeAssignment}, and they mentioned that in single-valued languages, the two are equivalent. Thus, in standard refinement type system which is single-valued, all type $\{\vnu : b ~|~ \phi \}$ can expressed as $\ctpers \ort{b}{\phi}$ in our system, which omit the different between $\ctcomma$ and $\ctstar$.

\paragraph{Another Persistent modality} Another persistent modality force make sure $\subExt \phi$ is persistent, which can be defined as: 
\begin{align*}
  r \models \chpers p &\text{ iff } \forall \sigma \in r. \{\sigma \} \models p
\end{align*}\noindent
It has a disadvantage, which still allowing the non-singleton resource to hold the property. Thus,
\begin{align*}
  &\chpers p \chland q \models \chpers p \chstar q \judgfail
  &\chpers p \models \chpers p \chstar \chpers p \judgfail
\end{align*}\noindent

